\section{GraphDrone.fit() Phase I-A Skeleton}

Phase I-A establishes the first experiment-independent package surface for GraphDrone. The implemented package introduces explicit configuration, descriptor validation, manifest-based portfolio loading, expert batching, token construction, bootstrap routing, and defer-to-FULL integration. This phase does not yet implement learned routing.

The key architectural decision is that benchmark runners no longer define the model boundary. Instead, the core object lives under \texttt{src/graphdrone\_fit/} and exposes a public \texttt{fit()} / \texttt{predict()} API. The benchmark layer is reduced to an adapter that converts registry-backed view declarations into typed descriptors.

For row $i$ and expert $v$, Phase I-A builds a token tensor containing:
\[
[\,\hat y_{i,v},\ \hat y_{i,v} - \hat y_{i,\mathrm{FULL}},\ \hat y_{i,v} - \bar y_i,\ q_{i,v},\ s_{i,v},\ d_v\,].
\]
Here $\bar y_i$ is the row-wise portfolio mean, used as a neutral within-row baseline for this phase only. Support encoding remains a zero-width placeholder tensor in Phase I-A; its explicit API keeps later contextual support work from leaking into benchmark scripts.

External review identified one real boundary issue: benchmark-specific family assignment logic was still partly hardcoded in the OpenML adapter. The final Phase I-A state fixes that by making family overrides explicit and by testing the override path directly, including \texttt{LOWRANK}.

Validation:
\begin{itemize}
\item \texttt{python -m py\_compile ...} passed
\item \texttt{pytest -q tests/graphdrone\_fit/test\_graphdrone\_fit.py tests/openml\_regression\_benchmark/test\_graphdrone\_fit\_adapter.py} passed with 6 tests
\end{itemize}

Phase I-A is therefore a valid checkpoint for the next implementation step: Phase I-B expert-factory work, still without learned contextual routing.
